
\section{Analisis Galat\label{sec:Runge-Kutta-Error-Analysis}\index{metode Runge-Kutta}}

Bagian \ref{sec:Runge-Kutta-Methods} berakhir dengan pengamatan yang
misterius (dan meresahkan?) bahwa PDB-Simpson tidak memenuhi harapan.
Berdasarkan pemecah PDB lainnya, kita memperkirakan laju kekonvergenan
PDB-Simpson sebesar $O(h^{4})$ karena Aturan Simpson, yang mendasari
PDB-Simpson, memiliki galat pemotongan lokal sebesar $O(h^{5})$.

Penjelasannya berakar pada fakta bahwa kita menyelesaikan PDB berbentuk
$\dot{y}=f(t,y)$, yang turunannya merupakan fungsi dari dua variabel,
$t$ dan $y$. Untuk memahami analisis galat, diperlukan penggunaan
intensif turunan parsial dan aturan rantai. Seperti biasa, kita menggunakan
teorema Taylor dan menuliskan
\[
y(t_{0}+h)=y(t_{0})+h\dot{y}(t_{0})+\frac{1}{2}h^{2}\ddot{y}(t_{0})+\frac{1}{6}h^{3}\dddot{y}(t_{0})+\cdots.
\]
Setiap turunan $y$ dapat diganti dengan suatu fungsi dari $f$ dan
turunan-turunan parsialnya, dimulai dengan $\dot{y}$, yang diberikan
oleh PDB yang hendak kita selesaikan.
\begin{eqnarray*}
\dot{y} & = & f(t,y)\\
\ddot{y}=\frac{d}{dt}\dot{y} & = & \frac{d}{dt}f(t,y)=f_{t}(t,y)+f_{y}(t,y)\dot{y}=f_{t}(t,y)+f_{y}(t,y)\cdot f(t,y)\\
 & \vdots
\end{eqnarray*}
Dengan menghilangkan penulisan eksplisit argumen $t$ dan $y$,
\begin{eqnarray*}
\dot{y} & = & f\\
\ddot{y} & = & f_{t}+f_{y}f\\
\dddot{y} & = & f_{tt}+f_{ty}f+(f_{yt}+f_{yy}f)f+f_{y}(f_{t}+f_{y}f)\\
 & = & f_{tt}+2f_{ty}f+f_{yy}f^{2}+f_{t}f_{y}+f_{y}^{2}f\\
 & \vdots
\end{eqnarray*}
sehingga $y(t_{0}+h)=y(t_{0})+h\dot{y}(t_{0})+\frac{1}{2}h^{2}\ddot{y}(t_{0})+\frac{1}{6}h^{3}\dddot{y}(t_{0})+\cdots$
dalam bentuk $f$ adalah
\[
y(t_{0}+h)=y(t_{0})+hf+\frac{1}{2}h^{2}(f_{t}+f_{y}f)+\frac{1}{6}h^{3}(f_{tt}+2f_{ty}f+f_{yy}f^{2}+f_{t}f_{y}+f_{y}^{2}f)+\cdots,
\]
dan sebagai pemecah PDB (dengan mengganti $y(t_{0})$ dengan $y_{i}$
dan $y(t_{0}+h)$ dengan $y_{i+1}$),
\begin{equation}
y_{i+1}=y_{i}+hf+\frac{1}{2}h^{2}(f_{t}+f_{y}f)+\frac{1}{6}h^{3}(f_{tt}+2f_{ty}f+f_{yy}f^{2}+f_{t}f_{y}+f_{y}^{2}f)+\cdots.\label{eq:twoVariableTaylorSeries}
\end{equation}
Menulis ulang polinomial Taylor berderajat tinggi dalam bentuk $f$
dengan cepat menjadi rumit. Kita akan berfokus pada analisis yang
hanya memerlukan $\dot{y}$, $\ddot{y}$, dan $\dddot{y}$.

Pemecah PDB pada Bagian \ref{sec:Runge-Kutta-Methods} berbentuk
\begin{eqnarray}
k_{1} & = & f(t_{i},y_{i})\nonumber \\
k_{2} & = & f\left(t_{i}+\beta_{2}h,y_{i}+\beta_{2}hk_{1}\right)\nonumber \\
k_{3} & = & f(t_{i}+\beta_{3}h,y_{i}+\beta_{3}hk_{2})\nonumber \\
 & \vdots\nonumber \\
k_{s} & = & f(t_{i}+\beta_{s}h,y_{i}+\beta_{s}hk_{s-1})\nonumber \\
y_{i+1} & = & y_{i}+h\left[\alpha_{1}k_{1}+\alpha_{2}k_{2}+\alpha_{3}k_{3}+\cdots+\alpha_{s}k_{s}\right].\label{eq:rungeKuttaGeneral}
\end{eqnarray}
Kita sebenarnya belum melihat pemecah PDB dengan $s>3$ di Bagian
\ref{sec:Runge-Kutta-Methods}, tetapi proses yang kita ikuti jelas
akan memerlukannya jika terdapat lebih dari tiga simpul dalam rumus
integrasi yang mendasarinya. 

Selisih antara $y(t_{0}+h)$ dari (\ref{eq:twoVariableTaylorSeries})
dan $y_{i+1}$ dari (\ref{eq:rungeKuttaGeneral}) adalah galat pemotongan
lokal pemecah PDB (galat ketika mengambil satu langkah). Untuk menuliskan
galat pemotongan ini dalam bentuk $O(h^{\ell})$, kita perlu mengembangkan
setiap $k_{j}$ ke dalam polinomial Taylor-nya. Diperlukan teorema
Taylor dalam dua variabel.
\begin{thm}
\label{thm:twoVariableTaylors}\index{Teorema!Taylor dua variabel}Misalkan
$f(t,y)$ beserta semua turunan parsialnya yang berorde $n+1$ atau
lebih rendah kontinu pada persegi panjang $D=\{(t,y):a\leq t\leq b,c\leq y\leq d\}$,
dan misalkan $(t_{0},y_{0})\in D$. Maka untuk setiap $(t,y)\in D$,
terdapat $\xi\in(a,b)$ dan $\mu\in(c,d)$ sedemikian sehingga 
\begin{eqnarray*}
f(t,y) & = & f(t_{0},y_{0})+\left[(t-t_{0})\cdot f_{t}(t_{0},y_{0})+(y-y_{0})\cdot f_{y}(t_{0},y_{0})\right]\\
 &  & +\frac{1}{2}\left[(t-t_{0})^{2}f_{tt}(t_{0},y_{0})+2(t-t_{0})(y-y_{0})\cdot f_{ty}(t_{0},y_{0})+(y-y_{0})^{2}f_{yy}(t_{0},y_{0})\right]\\
 &  & +\cdots+\\
 &  & \frac{1}{n!}\left[\sum_{j=0}^{n}\left(\begin{array}{c}
n\\
j
\end{array}\right)(t-t_{0})^{n-j}(y-y_{0})^{j}\frac{\partial^{n}f}{\partial t^{n-j}\partial y^{j}}(t_{0},y_{0})\right]\\
 &  & +\frac{1}{(n+1)!}\left[\sum_{j=0}^{n+1}\left(\begin{array}{c}
n+1\\
j
\end{array}\right)(t-t_{0})^{n+1-j}(y-y_{0})^{j}\frac{\partial^{n+1}f}{\partial t^{n+1-j}\partial y^{j}}(\xi,\mu)\right].
\end{eqnarray*}
\end{thm}

\noindent Seperti pada teorema Taylor (satu variabel), $n+1$ suku
pertama membentuk polinomial Taylor dan suku terakhir adalah suku
sisa. 

Sebagai ilustrasi, kita ambil $f(t,y)=-\frac{y}{t}+t^{2}$ dan menghitung
polinomial Taylor derajat dua beserta suku sisanya yang dikembangkan
di sekitar $(t_{0},y_{0})=(1,1)$. Untuk itu, kita memerlukan semua
turunan parsial $f$ hingga dan termasuk orde 3.
\begin{eqnarray*}
f_{t} & = & \frac{y}{t^{2}}+2t\\
f_{y} & = & -\frac{1}{t}\\
f_{tt} & = & -2\frac{y}{t^{3}}+2\\
f_{ty}=f_{yt} & = & \frac{1}{t^{2}}\\
f_{yy} & = & 0\\
f_{ttt} & = & 6\frac{y}{t^{4}}\\
f_{tty}=f_{tyt}=f_{ytt} & = & -\frac{2}{t^{3}}\\
f_{tyy}=f_{yty}=f_{yyt} & = & 0\\
f_{yyy} & = & 0.
\end{eqnarray*}
Dengan demikian,
\begin{eqnarray*}
f(1,1) & = & 0\\
f_{t}(1,1) & = & 3\\
f_{y}(1,1) & = & -1\\
f_{tt}(1,1) & = & 0\\
f_{ty}(1,1) & = & 1\\
f_{yy}(1,1) & = & 0\\
f_{ttt}(\xi,\mu) & = & 6\frac{\mu}{\xi^{4}}\\
f_{tty}(\xi,\mu) & = & -\frac{2}{\xi^{3}}\\
f_{tyy}(\xi,\mu) & = & 0\\
f_{yyy}(\xi,\mu) & = & 0.
\end{eqnarray*}
Oleh karena itu, polinomial Taylor derajat dua untuk $f(t,y)$ adalah
\begin{eqnarray*}
T_{2}(t,y) & = & f(1,1)+\left[(t-1)\cdot f_{t}(1,1)+(y-1)\cdot f_{y}(1,1)\right]\\
 &  & +\frac{1}{2}\left[(t-1)^{2}f_{tt}(1,1)+2(t-1)(y-1)\cdot f_{ty}(1,1)+(y-1)^{2}f_{yy}(1,1)\right]\\
 & = & 0+3(t-1)-(y-1)+0(t-1)^{2}+(t-1)(y-1)+0(y-1)^{2}\\
 & = & 3(t-1)-(y-1)+(t-1)(y-1)
\end{eqnarray*}
dengan suku sisa
\begin{eqnarray*}
R_{2}(t,y) & = & \frac{1}{6}\left[(t-1)^{3}f_{ttt}(\xi,\mu)+3(t-1)^{2}(y-1)f_{tty}(\xi,\mu)+3(t-1)(y-1)^{2}f_{tyy}(\xi,\mu)+(y-1)^{3}f_{yyy}(\xi,\mu)\right]\\
 & = & \frac{1}{6}\left[(t-1)^{3}\cdot6\frac{\mu}{\xi^{4}}-3(t-1)^{2}(y-1)\cdot\frac{2}{\xi^{3}}+3(t-1)(y-1)^{2}\cdot0+(y-1)^{3}\cdot0\right]\\
 & = & (t-1)^{3}\frac{\mu}{\xi^{4}}-(t-1)^{2}(y-1)\frac{1}{\xi^{3}}.
\end{eqnarray*}

Secara lebih umum, misalkan kita ingin memperoleh ekspansi polinomial
Taylor dari bentuk seperti $f(t_{i}+\beta_{j}h,y_{i}+\beta_{j}hk_{j-1})$,
sebagaimana dalam pemecah PDB kita. Dengan mengembangkannya di sekitar
$(t_{i},y_{i})$, kita tetapkan $t_{0}=t_{i}$, $y_{0}=y_{i}$, $t=t_{i}+\beta_{j}h$,
dan $y=y_{i}+\beta_{j}hk_{j-1}$. Jadi, $t-t_{0}=\beta_{j}h$ dan
$y-y_{0}=\beta_{j}hk_{j-1}$, sedangkan polinomial Taylor derajat
dua tanpa menuliskan argumen secara eksplisit $t_{i}$ dan $y_{i}$
pada ruas kanan adalah 
\begin{eqnarray*}
f(t_{i}+\beta_{j}h,y_{i}+\beta_{j}hk_{j-1}) & = & f+h\beta_{j}\left[f_{t}+k_{j-1}f_{y}\right]+\frac{1}{2}h^{2}\beta_{j}^{2}\left[f_{tt}+2k_{j-1}f_{ty}+k_{j-1}^{2}f_{yy}\right]
\end{eqnarray*}
dengan suku sisa $O(h^{3})$. 

Khususnya, ketika kita menetapkan $j=1$, $\beta_{j}=\beta_{1}=0$,
kita peroleh
\[
k_{1}=f(t_{i},y_{i})=f.
\]
Ketika kita menetapkan $j=2$, 
\begin{eqnarray*}
k_{2} & = & f\left(t_{i}+\beta_{2}h,y_{i}+\beta_{2}hk_{1}\right)\\
 & = & f+h\beta_{2}\left[f_{t}+ff_{y}\right]+\frac{1}{2}h^{2}\beta_{2}^{2}\left[f_{tt}+2ff_{ty}+f^{2}f_{yy}\right]+O(h^{3}).
\end{eqnarray*}
Perhitungan $k_{3}$ sedikit lebih rumit karena melibatkan $k_{2}^{2}$.
Namun, sebelum langsung menyelaminya, pertimbangkan apa yang akan
kita lakukan dengan $k_{3}$ terlebih dahulu. Setelah menghitung $k_{1}$,
$k_{2}$, dan $k_{3}$, kita akan menyubstitusikan masing-masing ke
dalam rumus
\begin{equation}
y_{i+1}=y_{i}+h\left[\alpha_{1}k_{1}+\alpha_{2}k_{2}+\alpha_{3}k_{3}\right]\label{eq:rungeKutta3ks}
\end{equation}
dan mengurangkan hasilnya dari (\ref{eq:twoVariableTaylorSeries}).
Untuk pembahasan ini, kita mencari metode dengan galat pemotongan
lokal sebesar $O(h^{4})$. Karena itu, kita hanya perlu mempertahankan
suku konstanta dan suku yang memuat faktor berupa $h^{3}$, $h^{2}$,
atau $h$ dalam persamaan (\ref{eq:rungeKutta3ks}). Suku-suku dengan
pangkat $h$ yang lebih tinggi tidak relevan. Suku-suku tersebut akan
terserap ke dalam $O(h^{4})$. Karena jumlah $\alpha_{1}k_{1}+\alpha_{2}k_{2}+\alpha_{3}k_{3}$
dikalikan dengan $h$, kita hanya perlu mempertahankan suku yang memuat
faktor hingga $h^{2}$ dalam $k_{1}$, $k_{2}$, dan $k_{3}$. Dengan
melihat ekspansi $k_{3}$:
\begin{eqnarray*}
k_{3} & = & f\left(t_{i}+\beta_{3}h,y_{i}+\beta_{3}hk_{2}\right)\\
 & = & f+h\beta_{3}\left[f_{t}+k_{2}f_{y}\right]+\frac{1}{2}h^{2}\beta_{3}^{2}\left[f_{tt}+2k_{2}f_{ty}+k_{2}^{2}f_{yy}\right]
\end{eqnarray*}
kita melihat bahwa hanya suku $\frac{1}{2}h^{2}\beta_{3}^{2}\cdot k_{2}^{2}f$
yang memuat $k_{2}^{2}$, dan suku itu sudah memiliki faktor $h^{2}$.
Akibatnya, kita hanya perlu menyertakan suku konstanta dari $k_{2}^{2}$.
Suku-suku lain dari $k_{2}^{2}$ menjadi bagian dari $O(h^{4})$.
Ternyata tidak terlalu buruk!
\begin{eqnarray*}
k_{2}^{2} & = & f^{2}+O(h).
\end{eqnarray*}
Demikian pula, ketika kita menyubstitusikan bentuk-bentuk untuk $k_{2}$
ke dalam $k_{3}$, kita akan berhati-hati menghindari suku apa pun
yang akan menghasilkan faktor $h$ berpangkat lebih besar dari 2:
\begin{eqnarray*}
k_{3} & = & f+h\beta_{3}\left[f_{t}+\left(f+h\beta_{2}\left[f_{t}+ff_{y}\right]\right)f_{y}\right]\\
 &  & +\frac{1}{2}h^{2}\beta_{3}^{2}\left[f_{tt}+2(f)f_{ty}+\left(f^{2}\right)f_{yy}\right]+O(h^{3})\\
 & = & f+h\beta_{3}f_{t}+h\beta_{3}ff_{y}+h^{2}\beta_{2}\beta_{3}(f_{t}f_{y}+ff_{y}^{2})\\
 &  & +\frac{1}{2}h^{2}\beta_{3}^{2}\left[f_{tt}+2ff_{ty}+f^{2}f_{yy}\right]+O(h^{3}).
\end{eqnarray*}

Setelah semua perhitungan terperinci itu, sekarang saat yang tepat
untuk bersandar sejenak dan melihat apa yang telah kita peroleh. Kita
telah mengembangkan semua suku dalam (\ref{eq:rungeKuttaGeneral})
untuk $s=3$ dan siap membandingkan hasilnya dengan ekspansi Taylor
PDB pada (\ref{eq:twoVariableTaylorSeries}). Selisih keduanya adalah
galat pemotongan lokal, sehingga kita akan mencari pangkat terkecil
dari $h$ yang tersisa setelah pengurangan. Dengan menyalin kedua
persamaan di sini agar mudah dirujuk, kita mengurangkan
\[
y_{i+1}=y_{i}+hf+\frac{1}{2}h^{2}(f_{t}+f_{y}f)+\frac{1}{6}h^{3}(f_{tt}+2f_{ty}f+f_{yy}f^{2}+f_{t}f_{y}+f_{y}^{2}f)+O(h^{4})
\]
dari
\begin{eqnarray*}
y_{i+1} & = & y_{i}+h\left[\alpha_{1}k_{1}+\alpha_{2}k_{2}+\alpha_{3}k_{3}\right]\\
 & = & y_{i}+h\alpha_{1}k_{1}+h\alpha_{2}k_{2}+h\alpha_{3}k_{3}\\
 & = & y_{i}+h\alpha_{1}f\\
 &  & +h\alpha_{2}\left(f+h\beta_{2}\left[f_{t}+ff_{y}\right]+\frac{1}{2}h^{2}\beta_{2}^{2}\left[f_{tt}+2ff_{ty}+f^{2}f_{yy}\right]+O(h^{3})\right)\\
 &  & +h\alpha_{3}\left(f+h\beta_{3}f_{t}+h\beta_{3}ff_{y}+h^{2}\beta_{2}\beta_{3}(f_{t}f_{y}+ff_{y}^{2})+\frac{1}{2}h^{2}\beta_{3}^{2}\left[f_{tt}+2ff_{ty}+f^{2}f_{yy}\right]+O(h^{3})\right).
\end{eqnarray*}
Suku konstanta (suku yang tidak memuat faktor $h$) pada setiap persamaan
hanyalah $y_{i}$, sehingga tidak ada konstanta yang tersisa setelah
pengurangan. Selisih suku-suku yang melibatkan $h$ adalah $hf-(h\alpha_{1}f+h\alpha_{2}f+h\alpha_{3}f)=hf(1-(\alpha_{1}+\alpha_{2}+\alpha))$,
jadi agar tidak ada $h$ yang tersisa dalam selisih, harus berlaku
\[
\alpha_{1}+\alpha_{2}+\alpha_{3}=1.
\]
Selisih suku-suku yang melibatkan $h^{2}f_{t}$ adalah $\frac{1}{2}h^{2}f_{t}-(h^{2}\alpha_{2}\beta_{2}f_{t}+h^{2}\alpha_{3}\beta_{3}f_{t})=h^{2}f_{t}(\frac{1}{2}-(\alpha_{2}\beta_{2}+\alpha_{3}\beta_{3}))$,
jadi agar tidak ada $h^{2}f_{t}$ yang tersisa dalam selisih, harus
berlaku
\[
\alpha_{2}\beta_{2}+\alpha_{3}\beta_{3}=\frac{1}{2}.
\]
Demikian pula, kita meninjau selisih suku-suku lainnya untuk memperoleh
syarat-syarat berikut bagi $\alpha_{j}$ dan $\beta_{j}$.

\begin{center}
\begin{tabular}{cc}
suku & menghasilkan syarat\tabularnewline
\hline 
$h^{2}f_{y}f$ & $\alpha_{2}\beta_{2}+\alpha_{3}\beta_{3}=\frac{1}{2}$\tabularnewline
$h^{3}f_{tt}$ & $\alpha_{2}\beta_{2}^{2}+\alpha_{3}\beta_{3}^{2}=\frac{1}{3}$\tabularnewline
$h^{3}f_{ty}f$ & $\alpha_{2}\beta_{2}^{2}+\alpha_{3}\beta_{3}^{2}=\frac{1}{3}$\tabularnewline
$h^{3}f_{yy}f^{2}$ & $\alpha_{2}\beta_{2}^{2}+\alpha_{3}\beta_{3}^{2}=\frac{1}{3}$\tabularnewline
$h^{3}f_{t}f_{y}$ & $\alpha_{3}\beta_{2}\beta_{3}=\frac{1}{6}$\tabularnewline
$h^{3}f_{y}^{2}f$ & $\alpha_{3}\beta_{2}\beta_{3}=\frac{1}{6}$\tabularnewline
\end{tabular}
\par\end{center}

\noindent Kita telah meninjau kedelapan suku yang berbeda, tetapi
hanya memperoleh empat syarat yang berlainan:
\begin{eqnarray}
\alpha_{1}+\alpha_{2}+\alpha_{3} & = & 1\nonumber \\
\alpha_{2}\beta_{2}+\alpha_{3}\beta_{3} & = & \frac{1}{2}\nonumber \\
\alpha_{2}\beta_{2}^{2}+\alpha_{3}\beta_{3}^{2} & = & \frac{1}{3}\nonumber \\
\alpha_{3}\beta_{2}\beta_{3} & = & \frac{1}{6}.\label{eq:rk3conditions}
\end{eqnarray}
Karena terdapat lima variabel dan hanya empat syarat, kita dapat menduga
bahwa terdapat beberapa pemecah PDB berbentuk (\ref{eq:rungeKuttaGeneral})
dengan $s=3$ dan galat pemotongan lokal $O(h^{4})$. 

Bukti dari Bagian \ref{sec:Runge-Kutta-Methods} menunjukkan bahwa
PDB-tertutup-terbuka seharusnya memiliki galat pemotongan lokal $O(h^{4}$).
Mari kita periksa. Untuk metode tersebut, kita memiliki
\[
\begin{gathered}\alpha_{1}=\frac{1}{4},\quad\alpha_{2}=0,\quad\alpha_{3}=\frac{3}{4}\\
\beta_{2}=\frac{1}{3},\quad\beta_{3}=\frac{2}{3},
\end{gathered}
\]
sehingga
\begin{eqnarray*}
\alpha_{1}+\alpha_{2}+\alpha_{3} & = & \frac{1}{4}+0+\frac{3}{4}=1\\
\alpha_{2}\beta_{2}+\alpha_{3}\beta_{3} & = & 0\cdot\frac{1}{3}+\frac{3}{4}\cdot\frac{2}{3}=\frac{1}{2}\\
\alpha_{2}\beta_{2}^{2}+\alpha_{3}\beta_{3}^{2} & = & 0\left(\frac{1}{3}\right)^{2}+\frac{3}{4}\left(\frac{2}{3}\right)^{2}=\frac{1}{3}\\
\alpha_{3}\beta_{2}\beta_{3} & = & \frac{3}{4}\cdot\frac{1}{3}\cdot\frac{2}{3}=\frac{1}{6}.
\end{eqnarray*}
Memang, PDB-tertutup-terbuka memenuhi semua syarat pemecah PDB dengan
galat pemotongan lokal (setidaknya) $O(h^{4})$. Sebenarnya, kita
harus menunjukkan bahwa setidaknya ada satu suku yang memuat $h^{4}$
tetap tersisa dalam selisih untuk membuktikan bahwa galat pemotongan
lokal tidak berderajat lebih tinggi.

Sebelum akhirnya menjawab apa yang terjadi pada PDB-Simpson, kerja
keras kita sejauh ini sudah cukup untuk memeriksa bahwa PDB-trapesium
dan PDB-terbuka memiliki galat pemotongan lokal $O(h^{3})$ dan bahwa
metode Euler memiliki galat pemotongan lokal $O(h^{2})$. Untuk PDB-trapesium,
kita memiliki $\alpha_{1}=\frac{1}{2}$, $\alpha_{2}=\frac{1}{2}$,
$\alpha_{3}=0$, $\beta_{2}=1$, dan $\beta_{3}$ tidak terdefinisi
(kita dapat menetapkan sembarang bilangan yang kita pilih karena $\alpha_{3}=0$
membuat $\beta_{3}$ tidak relevan bagi metode tersebut), sehingga
kita memperoleh
\begin{eqnarray*}
\alpha_{1}+\alpha_{2}+\alpha_{3} & = & \frac{1}{2}+\frac{1}{2}+0=1\\
\alpha_{2}\beta_{2}+\alpha_{3}\beta_{3} & = & \frac{1}{2}\cdot1+0=\frac{1}{2}\\
\alpha_{2}\beta_{2}^{2}+\alpha_{3}\beta_{3}^{2} & = & \frac{1}{2}\left(\frac{1}{3}\right)^{2}+0=\frac{1}{18}\neq\frac{1}{3}\\
\alpha_{3}\beta_{2}\beta_{3} & = & 0\neq\frac{1}{6}.
\end{eqnarray*}
Dua syarat pertama terpenuhi, tetapi dua syarat terakhir tidak. Namun,
ingat bahwa dua syarat pertama diturunkan dari suku $h$ dan $h^{2}$
sedangkan dua syarat terakhir diturunkan dari suku $h^{3}$. Jadi,
untuk PDB-trapesium, galat pemotongan lokalnya adalah $O(h^{3})$.

Untuk metode Euler, kita memiliki $\alpha_{1}=1$, $\alpha_{2}=\alpha_{3}=0$,
serta $\beta_{2}$ dan $\beta_{3}$ tidak terdefinisi (atau dapat
kita pilih sesuka hati), sehingga kita memperoleh
\begin{eqnarray*}
\alpha_{1}+\alpha_{2}+\alpha_{3} & = & 1+0+0=1\\
\alpha_{2}\beta_{2}+\alpha_{3}\beta_{3} & = & 0+0=0\neq\frac{1}{2}\\
\alpha_{2}\beta_{2}^{2}+\alpha_{3}\beta_{3}^{2} & = & 0+0=0\neq\frac{1}{3}\\
\alpha_{3}\beta_{2}\beta_{3} & = & 0\neq\frac{1}{6}.
\end{eqnarray*}
Persamaan kedua, yang diturunkan dari suku yang melibatkan $h^{2}$,
tidak terpenuhi, tetapi persamaan pertama, yang diturunkan dari suku
yang melibatkan $h$, terpenuhi, sehingga galat pemotongan lokal metode
Euler adalah $O(h^{2})$.

Terakhir, untuk PDB-Simpson, kita memiliki $\alpha_{1}=\frac{1}{6}$,
$\alpha_{2}=\frac{2}{3}$, $\alpha_{3}=\frac{1}{6}$, $\beta_{2}=\frac{1}{2}$,
dan $\beta_{3}=1$, sehingga kita memperoleh
\begin{eqnarray*}
\alpha_{1}+\alpha_{2}+\alpha_{3} & = & \frac{1}{6}+\frac{2}{3}+\frac{1}{6}=1\\
\alpha_{2}\beta_{2}+\alpha_{3}\beta_{3} & = & \frac{2}{3}\cdot\frac{1}{2}+\frac{1}{6}\cdot1=\frac{1}{2}\\
\alpha_{2}\beta_{2}^{2}+\alpha_{3}\beta_{3}^{2} & = & \frac{2}{3}\left(\frac{1}{2}\right)^{2}+\frac{1}{6}(1)^{2}=\frac{1}{3}\\
\alpha_{3}\beta_{2}\beta_{3} & = & \frac{1}{6}\cdot\frac{1}{2}\cdot1\neq\frac{1}{6}.
\end{eqnarray*}
Dua persamaan pertama terpenuhi, sehingga galat pemotongan lokalnya
(setidaknya) $O(h^{3})$, tetapi persamaan terakhir tidak terpenuhi,
sehingga galat pemotongan lokalnya tidak berorde lebih tinggi daripada
$O(h^{3})$. Tidak ada suku yang memuat faktor $h$ atau $h^{2}$
(yang tidak sekaligus memuat pangkat $h$ yang lebih tinggi) dalam
galat pemotongan lokal, tetapi suku $h^{3}\alpha_{3}\beta_{2}\beta_{3}(f_{t}f_{y}+ff_{y}^{2})=\frac{1}{6}h^{3}(f_{t}f_{y}+ff_{y}^{2})$
muncul, sehingga galatnya adalah $O(h^{3})$.

\subsection{Catatan tentang Konvensi dan Praktik}

Sejauh ini kita telah menurunkan lima pemecah PDB dengan hampir tidak
mengindahkan praktik yang mapan. Sudah saatnya kita memperbaiki hal
itu. Metode yang selama ini kita sebut PDB-trapesium (karena diturunkan
dari Aturan Trapesium) lebih dikenal sebagai metode Euler yang diperbaiki,\index{metode Euler yang diperbaiki}\index{metode!Euler yang diperbaiki|see {metode Euler yang diperbaiki}}
meskipun sebagian orang menyebutnya metode trapesium eksplisit.\index{metode trapesium eksplisit}\index{metode!trapesium eksplisit|see {metode trapesium eksplisit}}
Metode yang selama ini kita sebut PDB-tertutup-terbuka lebih dikenal
sebagai metode orde ketiga Heun.\index{metode orde ketiga Heun}\index{metode!orde ketiga Heun|see {metode orde ketiga Heun}}
Metode-metode ini mudah ditemukan dalam literatur. Keduanya merupakan
contoh prototipikal metode yang efisien. Metode Euler yang diperbaiki
memerlukan dua penghitungan nilai fungsi per langkah dan menghasilkan
galat pemotongan lokal $O(h^{3})$. Metode orde ketiga Heun memerlukan
tiga penghitungan nilai fungsi per langkah dan menghasilkan galat
pemotongan lokal $O(h^{4})$.

Metode yang selama ini kita sebut PDB-terbuka tidak diberi nama karena
tidak akan pernah digunakan dalam praktik. Metode ini tidak efisien
karena memerlukan tiga penghitungan nilai fungsi, tetapi hanya memiliki
galat pemotongan lokal $O(h^{3})$. Karena itu, kecil kemungkinan
Anda akan menemukannya dalam literatur karena metode ini tidak berguna
dalam praktik. Metode orde ketiga Heun maupun metode Euler yang diperbaiki
akan lebih baik daripada PDB-terbuka. Metode orde ketiga Heun menghasilkan
galat pemotongan yang lebih kecil dengan jumlah komputasi yang sama
(tiga penghitungan nilai fungsi), sedangkan metode Euler yang diperbaiki
menghasilkan galat pemotongan yang sama dengan komputasi lebih sedikit
(dua penghitungan nilai fungsi). PDB-Simpson memiliki kekurangan yang
sama dengan PDB-terbuka, sehingga kecil kemungkinan Anda menemukannya
dalam literatur. Metode ini juga tidak efisien.

Metode berbentuk (\ref{eq:rungeKuttaGeneral}) termasuk dalam kelas
metode yang disebut metode Runge-Kutta, dinamai menurut matematikawan
Jerman Carl Runge\index{Runge!Carl} dan Martin Kutta\index{Kutta!Martin}.
Gagasan dasar metode tersebut dikemukakan oleh Runge dalam sebuah
makalah yang diterbitkan pada tahun 1895; di sana Runge memperkenalkan
metode Euler yang diperbaiki dan metode lainnya. Karyanya dilanjutkan
oleh Heun\index{Heun!Karl}, yang melalui makalahnya pada tahun 1900
memperkenalkan metode orde ketiga Heun dan metode lainnya. Pada tahun
1901, Kutta menurunkan metode Runge-Kutta yang paling terkenal, yang
kini kadang-kadang disebut metode Runge-Kutta klasik atau metode Runge-Kutta
orde 4, RK4. Kita akan segera melihat bahwa metode ini merupakan modifikasi
dari PDB-Simpson.\cite{butcher}

\subsection{Metode Orde Lebih Tinggi}

Metode Runge-Kutta berorde lebih tinggi dapat diturunkan dengan mempertimbangkan
metode berbentuk (\ref{eq:rungeKuttaGeneral}) dengan sejumlah tahap,
$s>3$. Tentu saja, metode berorde lebih tinggi harus memenuhi lebih
banyak syarat. Faktanya, jumlah syarat bertambah lebih cepat ketika
orde yang diinginkan meningkat daripada pertambahan jumlah variabel
ketika jumlah tahap meningkat. Dengan kata lain, ada suatu titik ketika
jumlah tahap untuk mencapai orde $p$ melebihi $p$. Metode orde 1
dapat diturunkan dengan satu tahap (metode Euler) dan tidak dapat
kurang dari itu. Metode orde 2 dapat diturunkan dengan dua tahap (metode
Euler yang diperbaiki) dan tidak dapat kurang dari itu. Metode orde
3 dapat diturunkan dengan tiga tahap (metode orde ketiga Heun) dan
tidak dapat kurang dari itu. Metode orde 4 dapat diturunkan dengan
empat tahap (contohnya akan diberikan) dan tidak dapat kurang dari
itu. Namun, metode berorde $p$ dengan $p>4$ memerlukan jumlah tahap
$s>p$, yang pada gilirannya berarti lebih dari $p$ penghitungan
nilai fungsi. Jadi, metode yang paling efisien dijumpai pada orde
4 atau lebih rendah.

PDB-Simpson gagal mewujudkan potensinya karena tidak memiliki cukup
banyak tahap, bukan karena tidak ada rumus yang diturunkan dari Aturan
Simpson dengan galat pemotongan lokal $O(h^{5})$. Metode Runge-Kutta
klasik orde 4 (dengan galat pemotongan lokal $O(h^{5})$) memiliki
empat tahap dan diberikan oleh\index{metode RK4}\index{metode!RK4|see {metode RK4}}
\begin{eqnarray*}
k_{1} & = & f(t_{i},y_{i})\\
k_{2} & = & f\left(t_{i}+\frac{h}{2},y_{i}+\frac{h}{2}k_{1}\right)\\
k_{3} & = & f\left(t_{i}+\frac{h}{2},y_{i}+\frac{h}{2}k_{2}\right)\\
k_{4} & = & f\left(t_{i}+h,y_{i}+hk_{3}\right)\\
y_{i+1} & = & y_{i}+\frac{h}{6}\left[k_{1}+2k_{2}+2k_{3}+k_{4}\right].
\end{eqnarray*}
Bandingkan ini dengan PDB-Simpson:
\begin{eqnarray*}
k_{1} & = & f(t_{i},y_{i})\\
k_{2} & = & f\left(t_{i}+\frac{h}{2},y_{i}+\frac{h}{2}k_{1}\right)\\
k_{3} & = & f(t_{i+1},y_{i}+hk_{2})\\
y_{i+1} & = & y_{i}+\frac{h}{6}\left[k_{1}+4k_{2}+k_{3}\right].
\end{eqnarray*}
Keduanya sangat mirip. Jika tahap kedua PDB-Simpson dipisahkan menjadi
dua tahap, kita memperoleh metode Runge-Kutta orde 4. Itulah perbedaannya.
Dua tahap digunakan untuk menghampiri $\dot{y}(t_{i}+\frac{h}{2})$
alih-alih satu!

\noindent\begin{digression}{Penurunan Metode Runge-Kutta Orde 4 (Klasik)}\index{metode RK4}

\noindent Untuk menurunkan metode Runge-Kutta orde 4 mana pun, tahap-tahap
perhitungannya harus dikembangkan dalam polinomial Taylor derajat
tiga: 
\begin{eqnarray*}
f(t_{i}+\beta_{j}h,y_{i}+\beta_{j}hk_{j-1}) & = & f+h\beta_{j}\left[f_{t}+k_{j-1}f_{y}\right]+\frac{1}{2}h^{2}\beta_{j}^{2}\left[f_{tt}+2k_{j-1}f_{ty}+k_{j-1}^{2}f_{yy}\right]\\
 &  & +\frac{1}{6}h^{3}\beta_{j}^{3}\left[f_{ttt}+3k_{j-1}f_{tty}+3k_{j-1}^{2}f_{tyy}+k_{j-1}^{3}f_{yyy}\right]+O(h^{4})
\end{eqnarray*}
dan $f(t_{0},y_{0})$ harus dikembangkan dalam polinomial Taylor derajat
empat:
\[
y(t_{0}+h)=y(t_{0})+h\dot{y}(t_{0})+\frac{1}{2}h^{2}\ddot{y}(t_{0})+\frac{1}{6}h^{3}\dddot{y}(t_{0})+\frac{1}{24}h^{4}\ddddot{y}(t_{0})+O(h^{5}).
\]
Namun, $\ddddot{y}$, dalam bentuk $f$, adalah
\begin{eqnarray*}
\frac{d}{dt}(\dddot{y}) & = & \frac{d}{dt}\left(f_{tt}+2f_{ty}f+f_{yy}f^{2}+f_{t}f_{y}+f_{y}^{2}f\right)\\
 & = & f_{yyy}f^{3}+3f_{tyy}f^{2}+4f_{y}f_{yy}f^{2}+3f_{tty}f+5f_{ty}f_{y}f+f_{y}^{3}f\\
 &  & +3f_{t}f_{yy}f+f_{t}f_{y}^{2}+f_{tt}f_{y}+f_{ttt}+3f_{t}f_{ty}
\end{eqnarray*}
sehingga
\begin{eqnarray*}
y_{i+1} & = & y_{i}+hf+\frac{1}{2}h^{2}(f_{t}+f_{y}f)+\frac{1}{6}h^{3}(f_{tt}+2f_{ty}f+f_{yy}f^{2}+f_{t}f_{y}+f_{y}^{2}f)\\
 &  & +\frac{1}{24}h^{4}\left(f_{yyy}f^{3}+3f_{tyy}f^{2}+4f_{y}f_{yy}f^{2}+3f_{tty}f+5f_{ty}f_{y}f+f_{y}^{3}f\right.\\
 &  & \qquad\qquad\left.+3f_{t}f_{yy}f+f_{t}f_{y}^{2}+f_{tt}f_{y}+f_{ttt}+3f_{t}f_{ty}\right)+O(h^{5}).
\end{eqnarray*}

Selanjutnya,
\[
k_{1}=f(t_{i},y_{i})=f
\]
dan
\begin{eqnarray*}
k_{2} & = & f\left(t_{i}+\beta_{2}h,y_{i}+\beta_{2}hk_{1}\right)\\
 & = & f+h\beta_{2}\left[f_{t}+ff_{y}\right]+\frac{1}{2}h^{2}\beta_{2}^{2}\left[f_{tt}+2ff_{ty}+f^{2}f_{yy}\right]\\
 &  & +\frac{1}{6}h^{3}\beta_{2}^{3}\left[f_{ttt}+3ff_{tty}+3f^{2}f_{tyy}+f^{3}f_{yyy}\right]+O(h^{4}).
\end{eqnarray*}
Akibatnya, $k_{2}^{2}=f^{2}+2h\beta_{2}\left[f_{t}+ff_{y}\right]f+O(h^{2})$
dan $k_{2}^{3}=f^{3}+O(h)$. Oleh karena itu

\noindent
\begin{eqnarray*}
k_{3} & = & f+h\beta_{3}\left[f_{t}+k_{2}f_{y}\right]+\frac{1}{2}h^{2}\beta_{3}^{2}\left[f_{tt}+2k_{2}f_{ty}+k_{2}^{2}f_{yy}\right]\\
 &  & +\frac{1}{6}h^{3}\beta_{3}^{3}\left[f_{ttt}+3k_{2}f_{tty}+3k_{2}^{2}f_{tyy}+k_{2}^{3}f_{yyy}\right]\\
 & = & f+h\beta_{3}\left[f_{t}+\left(f+h\beta_{2}\left[f_{t}+ff_{y}\right]+\frac{1}{2}h^{2}\beta_{2}^{2}\left[f_{tt}+2ff_{ty}+f^{2}f_{yy}\right]\right)f_{y}\right]\\
 &  & +\frac{1}{2}h^{2}\beta_{3}^{2}\left[f_{tt}+2\left(f+h\beta_{2}\left[f_{t}+ff_{y}\right]\right)f_{ty}+\left(f^{2}+2h\beta_{2}\left[f_{t}+ff_{y}\right]f\right)f_{yy}\right]\\
 &  & +\frac{1}{6}h^{3}\beta_{3}^{3}\left[f_{ttt}+3ff_{tty}+3f^{2}f_{tyy}+f^{3}f_{yyy}\right]+O(h^{4})\\
 & = & f+h\beta_{3}\left[f_{t}+ff_{y}\right]+h^{2}\beta_{2}\beta_{3}\left[f_{t}+ff_{y}\right]f_{y}+\frac{1}{2}h^{2}\beta_{3}^{2}\left[f_{tt}+2ff_{ty}+f^{2}f_{yy}\right]\\
 &  & +\frac{1}{2}h^{3}\beta_{3}\beta_{2}^{2}\left[f_{tt}+2ff_{ty}+f^{2}f_{yy}\right]f_{y}+h^{3}\beta_{3}^{2}\beta_{2}\left[f_{t}+ff_{y}\right]\left[f_{ty}+ff_{yy}\right]\\
 &  & +\frac{1}{6}h^{3}\beta_{3}^{3}\left[f_{ttt}+3ff_{tty}+3f^{2}f_{tyy}+f^{3}f_{yyy}\right]+O(h^{4}).
\end{eqnarray*}
Jadi, $k_{3}^{2}=f^{2}+2h\beta_{3}\left[f_{t}+ff_{y}\right]f+O(h^{2})$
dan $k_{3}^{3}=f^{3}+O(h)$. Oleh karena itu
\begin{eqnarray*}
k_{4} & = & f+h\beta_{4}\left[f_{t}+k_{3}f_{y}\right]+\frac{1}{2}h^{2}\beta_{4}^{2}\left[f_{tt}+2k_{3}f_{ty}+k_{3}^{2}f_{yy}\right]\\
 &  & +\frac{1}{6}h^{3}\beta_{4}^{3}\left[f_{ttt}+3k_{3}f_{tty}+3k_{3}^{2}f_{tyy}+k_{3}^{3}f_{yyy}\right]+O(h^{4})\\
 & = & f+h\beta_{4}\left[f_{t}+\left(f+h\beta_{3}\left[f_{t}+ff_{y}\right]+h^{2}\beta_{2}\beta_{3}\left[f_{t}+ff_{y}\right]f_{y}+\frac{1}{2}h^{2}\beta_{3}^{2}\left[f_{tt}+2ff_{ty}+f^{2}f_{yy}\right]\right)f_{y}\right]\\
 &  & +\frac{1}{2}h^{2}\beta_{4}^{2}\left[f_{tt}+2\left(f+h\beta_{3}\left[f_{t}+ff_{y}\right]\right)f_{ty}+\left(f^{2}+2h\beta_{3}\left[f_{t}+ff_{y}\right]f\right)f_{yy}\right]\\
 &  & +\frac{1}{6}h^{3}\beta_{4}^{3}\left[f_{ttt}+3ff_{tty}+3f^{2}f_{tyy}+f^{3}f_{yyy}\right]+O(h^{4})\\
 & = & f+h\beta_{4}\left[f_{t}+ff_{y}\right]+h^{2}\beta_{3}\beta_{4}\left[f_{t}+ff_{y}\right]f_{y}+\frac{1}{2}h^{2}\beta_{4}^{2}\left[f_{tt}+2ff_{ty}+f^{2}f_{yy}\right]\\
 &  & +h^{3}\beta_{2}\beta_{3}\beta_{4}\left[f_{t}+ff_{y}\right]f_{y}^{2}+\frac{1}{2}h^{3}\beta_{4}\beta_{3}^{2}\left[f_{tt}+2ff_{ty}+f^{2}f_{yy}\right]f_{y}\\
 &  & +h^{3}\beta_{4}^{2}\beta_{3}\left[f_{t}+ff_{y}\right]\left[f_{ty}+ff_{yy}\right]+\frac{1}{6}h^{3}\beta_{4}^{3}\left[f_{ttt}+3ff_{tty}+3f^{2}f_{tyy}+f^{3}f_{yyy}\right]+O(h^{4}).
\end{eqnarray*}
Mencocokkan koefisien dalam 
\begin{eqnarray*}
y_{i+1} & = & y_{i}+hf+\frac{1}{2}h^{2}(f_{t}+f_{y}f)+\frac{1}{6}h^{3}(f_{tt}+2f_{ty}f+f_{yy}f^{2}+f_{t}f_{y}+f_{y}^{2}f)\\
 &  & +\frac{1}{24}h^{4}\left(f_{yyy}f^{3}+3f_{tyy}f^{2}+4f_{y}f_{yy}f^{2}+3f_{tty}f+5f_{ty}f_{y}f+f_{y}^{3}f\right.\\
 &  & \qquad\qquad\left.+3f_{t}f_{yy}f+f_{t}f_{y}^{2}+f_{tt}f_{y}+f_{ttt}+3f_{t}f_{ty}\right)+O(h^{5}).
\end{eqnarray*}
dengan koefisien dalam 
\[
y_{i+1}=y_{i}+h\left[\alpha_{1}k_{1}+\alpha_{2}k_{2}+\alpha_{3}k_{3}+\alpha_{4}k_{4}\right]
\]
hingga orde 4 menghasilkan syarat-syarat 

\noindent
\begin{eqnarray}
\alpha_{1}+\alpha_{2}+\alpha_{3}+\alpha_{4} & = & 1\label{eq:rk4cond1}\\
\alpha_{2}\beta_{2}+\alpha_{3}\beta_{3}+\alpha_{4}\beta_{4} & = & \frac{1}{2}\label{eq:rk4cond2}\\
\alpha_{2}\beta_{2}^{2}+\alpha_{3}\beta_{3}^{2}+\alpha_{4}\beta_{4}^{2} & = & \frac{1}{3}\label{eq:rk4cond3}\\
\alpha_{3}\beta_{2}\beta_{3}+\alpha_{4}\beta_{3}\beta_{4} & = & \frac{1}{6}\label{eq:rk4cond4}\\
\alpha_{2}\beta_{2}^{3}+\alpha_{3}\beta_{3}^{3}+\alpha_{4}\beta_{4}^{3} & = & \frac{1}{4}\label{eq:rk4cond5}\\
\alpha_{3}\beta_{3}^{2}\beta_{2}+\alpha_{4}\beta_{4}^{2}\beta_{3} & = & \frac{1}{8}\label{eq:rk4cond6}\\
2\alpha_{3}\beta_{3}^{2}\beta_{2}+2\alpha_{4}\beta_{4}^{2}\beta_{3}+\alpha_{3}\beta_{3}\beta_{2}^{2}+\alpha_{4}\beta_{4}\beta_{3}^{2} & = & \frac{1}{3}\label{eq:rk4cond7}\\
\alpha_{3}\beta_{3}^{2}\beta_{2}+\alpha_{4}\beta_{4}^{2}\beta_{3}+\alpha_{3}\beta_{3}\beta_{2}^{2}+\alpha_{4}\beta_{4}\beta_{3}^{2} & = & \frac{5}{24}\label{eq:rk4cond8}\\
\alpha_{3}\beta_{3}\beta_{2}^{2}+\alpha_{4}\beta_{4}\beta_{3}^{2} & = & \frac{1}{12}\label{eq:rk4cond9}\\
\alpha_{4}\beta_{2}\beta_{3}\beta_{4} & = & \frac{1}{24}.\label{eq:rk4cond10}
\end{eqnarray}
Setiap metode Runge-Kutta orde keempat dengan empat tahap ($s=4$)
yang berbentuk (\ref{eq:rungeKuttaGeneral}) harus memenuhi 10 persamaan
ini dengan hanya 7 derajat kebebasan (7 variabel). Persamaan-persamaan
tersebut membentuk himpunan dependen, atau hanya akan ada sedikit
solusi. Dalam upaya menyelesaikan sistem ini, kita selesaikan (\ref{eq:rk4cond10})
terhadap $\alpha_{4}$:
\[
\alpha_{4}=\frac{1}{24\beta_{2}\beta_{3}\beta_{4}}.
\]
Dengan menyubstitusikan rumus yang diperoleh untuk $\alpha_{4}$ ke
dalam (\ref{eq:rk4cond4}) dan menyelesaikannya terhadap $\alpha_{3}$:
\[
\alpha_{3}=\frac{4\beta_{2}-1}{24\beta_{2}^{2}\beta_{3}}.
\]
Dengan menyubstitusikan rumus yang diperoleh untuk $\alpha_{3}$ dan
$\alpha_{4}$ ke dalam (\ref{eq:rk4cond9}) dan menyelesaikannya terhadap
$\beta_{3}$:
\[
\beta_{3}=-4\beta_{2}^{2}+3\beta_{2}.
\]
Dengan menyubstitusikan rumus yang diperoleh untuk $\alpha_{3}$,
$\alpha_{4}$ dan $\beta_{3}$ ke dalam (\ref{eq:rk4cond6}) dan menyelesaikannya
terhadap $\beta_{4}$:
\[
\beta_{4}=(6-16\beta_{2}+16\beta_{2}^{2})\beta_{2}.
\]

\noindent Dengan menyubstitusikan rumus yang diperoleh untuk $\alpha_{3}$,
$\alpha_{4}$, $\beta_{3}$ dan $\beta_{4}$ ke dalam (\ref{eq:rk4cond2})
dan menyelesaikannya terhadap $\alpha_{2}$:
\[
\alpha_{2}=\frac{2-16\beta_{2}+52\beta_{2}^{2}-48\beta_{2}^{3}}{24\beta_{2}^{3}\left(3-4\beta_{2}\right)}.
\]
Dengan menyubstitusikan rumus yang diperoleh untuk $\alpha_{2}$,
$\alpha_{3}$, $\alpha_{4}$, $\beta_{3}$ dan $\beta_{4}$ ke dalam
(\ref{eq:rk4cond3}) dan menyederhanakannya:
\[
16\beta_{2}^{3}-12\beta_{2}^{2}+4\beta_{2}-1=0.
\]
Akar-akar persamaan terakhir ini adalah $\beta_{2}=\frac{1}{2},\frac{1\pm i\sqrt{7}}{8}$,
sehingga kita menyimpulkan bahwa $\beta_{2}=\frac{1}{2}$. Dengan
melakukan substitusi balik, kita peroleh
\begin{eqnarray*}
\beta_{2} & = & \frac{1}{2}\\
\alpha_{2} & = & \frac{1}{3}\\
\beta_{4} & = & 1\\
\beta_{3} & = & \frac{1}{2}\\
\alpha_{3} & = & \frac{1}{3}\\
\alpha_{4} & = & \frac{1}{6}.
\end{eqnarray*}
Dengan menyubstitusikan nilai $\alpha_{2}$, $\alpha_{3}$, dan $\alpha_{4}$
ke dalam (\ref{eq:rk4cond1}), kita peroleh
\[
\alpha_{1}=\frac{1}{6}.
\]
Ketujuh nilai ini adalah solusi riil simultan yang unik dari persamaan
(\ref{eq:rk4cond10}), (\ref{eq:rk4cond4}), (\ref{eq:rk4cond9}),
(\ref{eq:rk4cond6}), (\ref{eq:rk4cond2}), (\ref{eq:rk4cond3}),
dan (\ref{eq:rk4cond1}). Jadi, ketujuh parameter ditentukan oleh
tujuh dari sepuluh syarat tersebut. Tinggal ditunjukkan bahwa ketujuh
nilai ini juga memenuhi (\ref{eq:rk4cond5}), (\ref{eq:rk4cond7}),
dan (\ref{eq:rk4cond8}), dan memang demikian. Terakhir, perhatikan
bahwa inilah nilai-nilai parameter untuk metode Runge-Kutta (klasik)
orde 4.

\noindent\end{digression}

\subsection{Konsep Kunci}
\begin{description}
\item [{Teorema~Taylor~dalam~dua~variabel:}] \index{Teorema!Taylor dua variabel}Misalkan
$f(t,y)$ beserta semua turunan parsialnya yang berorde $n+1$ atau
lebih rendah kontinu pada persegi panjang $D=\{(t,y):a\leq t\leq b,c\leq y\leq d\}$,
dan misalkan $(t_{0},y_{0})\in D$. Maka untuk setiap $(t,y)\in D$,
terdapat $\xi\in(a,b)$ dan $\mu\in(c,d)$ sedemikian sehingga 
\begin{eqnarray*}
f(t,y) & = & f(t_{0},y_{0})+\left[(t-t_{0})\cdot f_{t}(t_{0},y_{0})+(y-y_{0})\cdot f_{y}(t_{0},y_{0})\right]\\
 &  & +\frac{1}{2}\left[(t-t_{0})^{2}f_{tt}(t_{0},y_{0})+2(t-t_{0})(y-y_{0})\cdot f_{ty}(t_{0},y_{0})+(y-y_{0})^{2}f_{yy}(t_{0},y_{0})\right]\\
 &  & +\cdots+\\
 &  & \frac{1}{n!}\left[\sum_{j=0}^{n}\left(\begin{array}{c}
n\\
j
\end{array}\right)(t-t_{0})^{n-j}(y-y_{0})^{j}\frac{\partial^{n}f}{\partial t^{n-j}\partial y^{j}}(t_{0},y_{0})\right]\\
 &  & +\frac{1}{(n+1)!}\left[\sum_{j=0}^{n+1}\left(\begin{array}{c}
n+1\\
j
\end{array}\right)(t-t_{0})^{n+1-j}(y-y_{0})^{j}\frac{\partial^{n+1}f}{\partial t^{n+1-j}\partial y^{j}}(\xi,\mu)\right].
\end{eqnarray*}
\end{description}
\noindent\startexercises 

\subsection{Latihan}
\begin{enumerate}
\item \label{exc:ode-errorAnalysis}Tentukan secara analitis galat pemotongan
lokal untuk pemecah PDB yang diturunkan dalam latihan \vref{exc:ode-rungeKutta}.
Bandingkan hasilnya dengan galat pemotongan lokal rumus integrasi
yang mendasarinya. Apakah keduanya sama? Bandingkan pula dengan laju
kekonvergenan yang ditentukan secara eksperimental (lihat latihan
\vref{exc:ode-rungeKutta-16}). Apakah nilainya satu derajat lebih
tinggi, seperti yang seharusnya diharapkan?  \haspartwithsolution  \haspartwithanswer 
\item \label{exc:ode-errorAnalysis-1}Lakukan satu langkah metode Runge-Kutta
orde empat untuk menyelesaikan $\dot{y}=ty$ dengan $y(1)=0.5$ dan
$h=1$, sehingga menghampiri $y(2)$. Bandingkan jawaban Anda dengan
jawaban pada Bagian \ref{sec:TaylorMethods} latihan \ref{exc:ode-taylor}c
\vpageref{exc:ode-taylor} tempat Anda menggunakan metode Euler dengan
dua langkah. Solusi eksaknya adalah $y(2)=\frac{e^{3/2}}{2}\approx2.240844535169032$. \hasasolution 
\item Jelaskan metode Euler yang diperbaiki secara geometris dan dengan
kata-kata Anda sendiri.
\item \octave\ \label{exc:ode-errorAnalysis-2}Tulislah fungsi Octave yang
mengimplementasikan metode Euler yang diperbaiki (sama seperti latihan
\vref{exc:ode-rungeKutta-12}, tetapi kali ini metode tersebut memiliki
nama yang tepat). \hasananswer 
\item \octave\ \label{exc:ode-errorAnalysis-3}Tulislah fungsi Octave yang
mengimplementasikan metode orde ketiga Heun (sama seperti latihan
\vref{exc:ode-rungeKutta-13}, tetapi kali ini metode tersebut memiliki
nama yang tepat). \hasananswer 
\item \octave\ \label{exc:ode-errorAnalysis-4}Tulislah fungsi Octave yang
mengimplementasikan metode RK4. \hasananswer 
\item \octave\ Gunakan kode Anda dari latihan \ref{exc:ode-errorAnalysis-4}
untuk menghitung $y(2)$ untuk PDB dalam latihan \ref{exc:ode-taylor}
\vpageref{exc:ode-taylor} dengan ukuran langkah $h=0.05$.  \haspartwithsolution  \haspartwithanswer 
\end{enumerate}
\noindent\finishexercises 
